\graphicspath{{chapters/03related/graphics}}

\chapter{Related work}

%This chapter covers relevant (and typically, recent) research
%which you build upon (or improve upon). There are two complementary
%goals for this chapter:
%\begin{enumerate}
%  \item to show that you know and understand the state of the art; and
%  \item to put your work in context
%\end{enumerate}
%
%Ideally you can tackle both together by providing a critique of
%related work, and describing what is insufficient (and how you do
%better!)
%
%The related work chapter should usually come either near the front or
%near the back of the dissertation. The advantage of the former is that
%you get to build the argument for why your work is important before
%presenting your solution(s) in later chapters; the advantage of the
%latter is that don't have to forward reference to your solution too
%much. The correct choice will depend on what you're writing up, and
%your own personal preference.

%In the related work part, briefly talk about how existing works differ from yours. For instance, DistServe and Splitwise are used for efficient LLM inference and serving, but they fail to account for handling hardware heterogeneity and might obtain poor performance when deployed on a fully heterogeneous cluster with various GPU resources and network bandwidth. Additionally, HexGen and ThunderServe are responsible for LLM serving with heterogeneous clusters, but they fail to consider the unique characteristics of MoE architecture ……

% Existing Schedulers
\section{Case Studies}

\todo{

(Already complete. Notes in Obsidian, need to be transferred)

\subsection{DistServe}
DistServe \autocite{zhongDistServeDisaggregatingPrefill2024} (Introduce Prefill/Decode).

\subsection{Splitwise}
Splitwise \autocite{patelSplitwiseEfficientGenerative2024} (Introduce Prefill/Decode).



\subsection{Thunderserve}
Thunderserve \cite{jiangThunderServeHighperformanceCostefficient2025}

\subsection{HexGen}
HexGen \autocite{jiangHexGenGenerativeInference2024}


\subsection{Summary}

HexGen and ThunderServe are responsible for LLM serving with heterogeneous clusters, but they fail to consider the unique characteristics of MoE architecture.

\section{Future}

Possible Improvements, Lack of Moe

}